\documentclass[11pt,a4paper]{ctexart}
\input{/Users/shaoyikai/Desktop/考研/22408-notes-skill/assets/note-preamble.tex}
\renewcommand{\NoteFooter}{22408 一页纸笔记 · 数学二}
\linespread{1.14}
\begin{document}
\notetitle{正交变换下标准型的换序}{数学二}{线性代数·二次型正交变换}
\sechead{一、核心原理}

\notebox{设 $P$ 为正交阵，在 $x=Py$ 下 $f=x^{\mathrm T}Ax$ 的标准型为 $\lambda_1y_1^2+\cdots+\lambda_ny_n^2$，则 $AP=P\Lambda$，$\Lambda=\mathrm{diag}(\lambda_1,\dots,\lambda_n)$。}

即 \textbf{$P$ 的第 $i$ 列是 $\lambda_i$ 对应的特征向量}，标准型第 $i$ 个系数就是 $P$ 第 $i$ 列所配的特征值。

\notebox{推论：换用另一正交阵 $Q$，\textbf{特征值集合不变，只是随列序重新排列}。}

\sechead{二、本题解法（三步）}

\textbf{题}：$f(x_1,x_2,x_3)$ 在 $x=Py$ 下标准型为 $2y_1^2+y_2^2-y_3^2$，$P=(e_1,e_2,e_3)$。若 $Q=(e_1,-e_3,e_2)$，求 $x=Qy$ 下的标准型。

\textbf{第 1 步} 标准型系数即特征值，与 $P$ 的列逐位配对：

\[
Ae_1=2e_1,\qquad Ae_2=e_2,\qquad Ae_3=-e_3
\]

\textbf{第 2 步} 把 $Q$ 的每列翻译成特征值：

\begin{center}\small
\begin{tabularx}{\linewidth}{@{}>{\raggedright\arraybackslash}X>{\raggedright\arraybackslash}X>{\raggedright\arraybackslash}X@{}}
\toprule
$Q$ 的列 & 计算 & 特征值 \\
\midrule
$e_1$ & $Ae_1=2e_1$ & $2$ \\
$-e_3$ & $A(-e_3)=-Ae_3=-1\cdot(-e_3)$ & $-1$ \\
$e_2$ & $Ae_2=e_2$ & $1$ \\
\bottomrule
\end{tabularx}
\end{center}

\textbf{第 3 步} 按 $Q$ 的列序写标准型：

\[
\boxed{\,f=2y_1^2-y_2^2+y_3^2\,}
\]

\textbf{答案：(A)}

\warnline{关键陷阱：$Q$ 第 2 列是 $-e_3$，而 $e_3$ 配特征值 $-1$，故系数仍为 $-1$——\textbf{特征向量取负号不改变特征值}。}

\sechead{三、秒杀口诀}

\notebox{\textbf{"看列不看值，列是谁就配谁的特征值。"}}

\[
\underbrace{(e_1,\ e_2,\ e_3)}_{(2,\ 1,\ -1)}
\ \xrightarrow{\ \text{换成}\ Q\ }\
\underbrace{(e_1,\ -e_3,\ e_2)}_{(2,\ -1,\ 1)}
\]

即直接在系数序列 $(2,1,-1)$ 上做与列相同的位置置换，得到 $(2,-1,1)$。

\sechead{四、易错点}

\warnline{别把 $Q$ 的列当新特征向量重求：正交变换不改变 $A$，特征值恒为 $2,1,-1$。}

\warnline{别看到 $-e_3$ 就以为系数变号：$-\alpha$ 与 $\alpha$ 同属一个特征值。}

\warnline{分清"系数"与"分量"：配对关系由 $AP=P\Lambda$ 决定，系数是数、列是向量。}

\warnline{正交变换下\textbf{惯性指数不变}：正 $2$ 负 $1$。(B)(D) 为正 $3$ 负 $0$，可先行排除。}

\par\addvspace{3pt}\noindent{\color{line}\hrule height 0.4pt}\par\addvspace{3pt}

\emph{依据：考研数学二线性代数 · 二次型与正交变换。}

\end{document}
